\section{Conclusions}\label{sec:conclusion}

We have proposed a method that makes the diagnosability of failure modes a design-stage criterion for choosing between architectural alternatives and, jointly, for choosing the observation suite that will monitor the chosen architecture. Its core is a scale-free, self-calibrated statistic, the signal-to-drift ratio, from which a minimum detectable extent and a detectable share can be computed for every combination of alternative, suite and failure-mode class, and a six-step procedure that turns these into a co-selection and a transferability specification.

Applied to three wind-turbine generator topologies subjected to identical winding short-circuits on open benchmark data, the method showed that the same design-level fault is not the same physical fault across architectures, that the control structure decides where in the drive a fault becomes visible, that the cheapest observation suite meeting a given detectability requirement is architecture-specific---terminal voltage transducers are decisive for the PMSG and superfluous for the SCIG---and that a learned detector transfers across architectures only when its features are referenced to the monitored unit's own healthy behaviour. The findings are robust to the method's tuning choices and are reported with recording-level uncertainty; they are limited by one physical unit per alternative, laboratory scale and a single failure-mode class.

Three extensions follow directly. First, the design-stage instantiation: simulated populations of each alternative under parameter variation, with the physical benches as validation points, would turn the demonstration's machine-level differences into population-level statements and would let the method be applied before hardware exists. Second, other failure-mode classes and observables---bearing and rotor faults with vibration and flux observables---would test the method's claim to generality within the same domain, and a mechanical or electrochemical family would test it across domains. Third, the cost side of the selection, treated here as an ordinal ranking of suites, deserves a proper model that includes integration, maintenance and the cost of missed detection, so that the co-selection can be posed as an optimisation. The code released with this paper reproduces every result from the public data and is intended as the starting point for those extensions.
